\section{Real Project Case Study}

\subsection{System Description}
The demonstration project consists of a minimal shopping application with a
\textbf{FastAPI}\footnote{Only open-source packages are used.} backend and a
React front end. The backend exposes REST endpoints for login, logout, browsing
products, manipulating the cart and carrying out the checkout flow. Each
endpoint returns JSON and stores session data in memory. A simplified selection
of products is provided so that the state transitions can be exercised without a
persistent database. Listing\,\ref{lst:machine_snippet} shows how the business
flow is encoded in the accompanying \texttt{state.machine} file.

The frontend is implemented with React and uses a context provider to maintain
its finite-state machine (FSM). States such as\texttt{Login},
\texttt{Browsing},\texttt{CartEditing} and \texttt{Checkout} are enumerated in
\texttt{AppContext.jsx}. Actions dispatched from UI components cause state
changes and trigger API calls through a small service layer. Pages in
\texttt{src/pages} map directly to these FSM states so that the current view is
always consistent with the application logic.

\begin{lstlisting}[caption={Excerpt of shopping state machine}, label={lst:machine_snippet}]
Login
  [Actions]
    login with credentials    ==>  Browsing      when  ${EMAIL} == ${VALID_EMAIL}  and  ${PASSWORD} == ${VALID_PASSWORD}
    login with credentials    ==>  LoginFailed   otherwise

LoginFailed
  [Actions]
    try login again          ==>  Login
    exit application         ==>  SessionEnded

Browsing
  [Actions]
    add product to cart       ==>  CartEditing    args    ${PRODUCT_ID}    ${QUANTITY}
    logout                    ==>  SessionEnded
\end{lstlisting}

Figure\,\ref{fig:shopping_fsm} illustrates the overall FSM.

\begin{figure}[ht]
    \centering
    \includegraphics[width=\linewidth]{shopping-fsm.png}
    \caption{State machine driving the shopping application.}
    \label{fig:shopping_fsm}
\end{figure}

\subsection{Typical Workflow in Practice}
An end user begins on the login page where credentials are submitted to
\texttt{/login}. On success the server issues a session token and the FSM moves
to the \texttt{Browsing} state. Product data retrieved from \texttt{/products}
populate the catalogue view. The user may add items to the cart, triggering the
\texttt{/cart/add} endpoint and a transition to \texttt{CartEditing}. While
editing the cart, quantities can be adjusted or items removed via
\texttt{/cart/update} and \texttt{/cart/remove}. Selecting ``Proceed to
Checkout'' calls \texttt{/checkout/start} and moves the FSM to \texttt{Checkout}.

During checkout a form captures shipping details. Payment is simulated through
\texttt{/checkout/pay} where success or failure is controlled by a flag. If the
payment succeeds the next state is \texttt{OrderConfirmed}. The frontend
displays an order summary along with a newly generated order ID. A payment error
leads instead to \texttt{PaymentFailed}. In this state the user can retry the
payment, cancel the order (returning to \texttt{CartEditing}) or logout. The
session ultimately ends through the \texttt{/logout} endpoint which clears
server-side data and returns to the initial \texttt{Login} state.

The workflow therefore mirrors the state machine: every API response results in
a deterministic state change on the client. The small number of states keeps the
model manageable while still covering typical e-commerce behaviour such as
cancelled checkouts or invalid login attempts.

\subsection{Different Test Generation Results}
The \texttt{state.machine} file feeds directly into the generator from
\texttt{tests/generator}. Three strategies are provided: a
\textbf{Depth-First Search} (DFS) strategy, a \textbf{Random} strategy and an
\textbf{All-Pairs} strategy. Each produces Robot Framework suites containing a
set of parametrised tests with coverage information automatically inserted at
the top of each generated file.

Table~\ref{tab:shopping_coverage_summary} compares the resulting test suites,
showing how different strategies achieve varying levels of coverage with
different numbers of test cases.

\begin{table}[ht]
\centering
\caption{Comparison of test generation strategies for the shopping application}
\label{tab:shopping_coverage_summary}
\begin{tabular}{|l|c|c|c|}
\hline
\textbf{Strategy} & \textbf{Test Cases} & \textbf{States Covered} & \textbf{Actions Covered} \\
\hline
Depth-First & 52 & 6/8 (75\%) & 10/17 (59\%) \\
\hline
Random & 52 & 8/8 (100\%) & 17/17 (100\%) \\
\hline
All-Pairs & 11 & 7/8 (88\%) & 12/17 (71\%) \\
\hline
\end{tabular}
\end{table}

Table~\ref{tab:shopping_detailed_coverage} provides detailed breakdown of which
specific states and actions are covered by each strategy in the shopping
application.

\begin{table}[ht]
\centering
\caption{Detailed coverage breakdown for shopping application strategies}
\label{tab:shopping_detailed_coverage}
\small
\begin{tabular}{|l|p{8cm}|p{3cm}|}
\hline
\textbf{Strategy} & \textbf{Covered States} & \textbf{Uncovered States} \\
\hline
Depth-First & 
\begin{tabular}[c]{@{}l@{}}
Browsing, CartEditing, \\
Checkout, Login, \\
OrderConfirmed, SessionEnded
\end{tabular} & LoginFailed, PaymentFailed \\
\hline
Random & All states (8/8) & None \\
\hline
All-Pairs & 
\begin{tabular}[c]{@{}l@{}}
Browsing, CartEditing, Checkout, \\
Login, LoginFailed, \\
OrderConfirmed, SessionEnded
\end{tabular} & PaymentFailed \\
\hline
\end{tabular}

\vspace{0.5cm}

\begin{tabular}{|l|p{9cm}|p{3cm}|}
\hline
\textbf{Strategy} & \textbf{Key Uncovered Actions} & \textbf{Missing Scenarios} \\
\hline
Depth-First & 
\begin{tabular}[c]{@{}l@{}}
• login with credentials (Login $\to$ LoginFailed) \\
• process payment (Checkout $\to$ PaymentFailed) \\
• cancel order (PaymentFailed $\to$ CartEditing) \\
• retry payment (PaymentFailed $\to$ Checkout)
\end{tabular} & 
\begin{tabular}[c]{@{}c@{}}
Login failures \\
Payment failures
\end{tabular} \\
\hline
Random & None (full coverage) & None \\
\hline
All-Pairs & 
\begin{tabular}[c]{@{}l@{}}
• process payment (Checkout $\to$ PaymentFailed) \\
• cancel order (PaymentFailed $\to$ CartEditing) \\
• retry payment (PaymentFailed $\to$ Checkout) \\
• logout (PaymentFailed $\to$ SessionEnded)
\end{tabular} & Payment failures \\
\hline
\end{tabular}
\end{table}

The random strategy achieves complete coverage because it explores a wider range
of variable assignments. DFS follows deterministic paths and therefore misses
some error handling states. All-Pairs trades thoroughness for brevity by
exercising only the combinations necessary to cover each pair of input
parameters at least once.

\subsection{Discussion and Future Work}
This small project demonstrates how a concise FSM can drive both the frontend
logic and the automated tests. Expressing the workflow in a separate
\texttt{.machine} file keeps the model readable and enables multiple generation
strategies without rewriting the test logic. While the random strategy yields the
best coverage, it also introduces redundant paths and longer execution time. DFS
is predictable but overlooks negative scenarios unless depth limits are raised.
All-Pairs generates a compact suite, yet it may miss complex interactions beyond
pairwise combinations.

Future enhancements could incorporate heuristics to prioritise rarely exercised
transitions or to balance coverage against runtime. The current generator emits
plain Robot Framework scripts; integrating screenshots or API traces would aid
debugging. Scaling the approach to larger systems would require grouping states
into higher-level regions and supporting parameterised submachines. Extending the
frontend to log state transitions could also help correlate UI behaviour with the
generated tests. Overall, the case study highlights both the practicality of
model-based testing and the importance of choosing a generation strategy that
suits the target system.

\subsection{Parameter Impact Analysis}
To understand how different parameter configurations affect test generation
effectiveness, we conducted a comprehensive analysis varying the \texttt{max\_tests}
and \texttt{max\_actions} parameters across all three strategies. This analysis
provides insights into the trade-offs between test suite size, execution time,
and coverage completeness.

\subsubsection{Coverage Performance by Strategy}
Figure~\ref{fig:coverage_overview} presents a comprehensive comparison of coverage
performance across the three test generation strategies. The average coverage
chart reveals that the Random strategy consistently outperforms both Depth-First
and All-Pairs strategies, achieving approximately 93\% state coverage and 82\%
action coverage on average. The Depth-First strategy shows more conservative
performance with 83\% state coverage and 66\% action coverage, while All-Pairs
demonstrates moderate effectiveness with 78\% state coverage and 62\% action
coverage.

\begin{figure}[ht]
    \centering
    \includegraphics[width=\linewidth]{coverage_analysis_overview.png}
    \caption{Coverage analysis overview showing strategy performance and parameter effects.}
    \label{fig:coverage_overview}
\end{figure}

The coverage distribution visualization using violin plots reveals important
differences in strategy behaviour. The Random strategy exhibits a wide
distribution with high mean coverage (approximately 87\%), indicating consistent
performance across different parameter combinations but with some variability.
The Depth-First strategy shows a narrower distribution centred around 75\%,
suggesting more predictable but limited coverage. The All-Pairs strategy
demonstrates bimodal behaviour with peaks around 65\% and 95\%, reflecting its
sensitivity to parameter configuration.

The parameter effect analysis demonstrates that increasing \texttt{max\_tests}
has a profound impact on the Random strategy, with coverage improving from 68\%
to 98\% as the test limit increases from 10 to 200 tests. Conversely, the
Depth-First strategy shows minimal improvement (67\% to 74\%), suggesting that
additional tests provide diminishing returns when following deterministic paths.
The All-Pairs strategy remains relatively stable around 74\%, as expected from
its combinatorial nature.

The \texttt{max\_actions} parameter shows different patterns. The Random strategy
benefits significantly from longer test sequences, improving from 57\% to 92\%
coverage as \texttt{max\_actions} increases from 5 to 30. The Depth-First
strategy shows gradual improvement from 83\% to 67\%, with an interesting
plateau effect at higher values. The All-Pairs strategy demonstrates consistent
improvement from 77\% to 81\%, indicating moderate sensitivity to sequence length.

\subsubsection{Detailed Distribution Analysis}
Figure~\ref{fig:coverage_distribution} provides deeper insights into the coverage
distribution patterns through histogram analysis and statistical summaries. The
histogram reveals that the Random strategy achieves the most consistent high
coverage, with most test configurations falling between 80-100\% coverage. The
Depth-First strategy shows a narrower distribution with most results clustered
around 70-75\%, while the All-Pairs strategy exhibits a bimodal distribution
with distinct peaks at approximately 65\% and 95\%.

\begin{figure}[ht]
    \centering
    \includegraphics[width=\linewidth]{coverage_distribution_detailed.png}
    \caption{Detailed coverage distribution analysis with histograms and statistical summaries.}
    \label{fig:coverage_distribution}
\end{figure}

The statistical summary with error bars confirms these observations, showing that
the Random strategy achieves the highest mean coverage with moderate standard
deviation, indicating reliable performance. The Depth-First strategy demonstrates
low variability but also lower mean coverage, while the All-Pairs strategy shows
high variability, reflecting its sensitivity to parameter selection.

\subsubsection{Parameter Sensitivity Tables}
Tables~\ref{tab:coverage_statistics} and \ref{tab:best_configurations} provide
quantitative analysis of strategy performance and optimal parameter combinations.
The coverage statistics reveal that while the Random strategy achieves the
highest mean coverage (87.3\%), it also exhibits the highest variability
(standard deviation of 13.2\%). This suggests that while Random can achieve
excellent coverage, careful parameter tuning is essential for consistent results.

\begin{table}[h!]
\centering
\caption{Coverage Statistics by Strategy}
\label{tab:coverage_statistics}
\begin{tabular}{|l|c|c|c|c|c|c|}
\hline
\textbf{Strategy} & \multicolumn{2}{c|}{\textbf{State Coverage (\%)}} & \multicolumn{2}{c|}{\textbf{Action Coverage (\%)}} & \multicolumn{2}{c|}{\textbf{Total Coverage (\%)}} \\
\cline{2-7}
 & Mean & Std & Mean & Std & Mean & Std \\
\hline
DepthFirst & 83.0 & 7.4 & 66.5 & 4.1 & 74.8 & 4.9 \\
\hline
Random & 93.2 & 10.1 & 81.5 & 16.3 & 87.3 & 13.2 \\
\hline
\end{tabular}
\end{table}

The Depth-First strategy shows the most consistent performance with low standard
deviation (4.9\%), making it predictable for scenarios where reliability is
prioritised over maximum coverage. The absence of All-Pairs results in the
statistics table indicates that this strategy encountered compatibility issues
with the rule-based state machine, highlighting a limitation when applying
combinatorial testing to systems with complex constraints.

Table~\ref{tab:best_configurations} identifies the optimal parameter combinations
for achieving maximum coverage. Notably, all top-performing configurations use
high values for both \texttt{max\_tests} (100-200) and \texttt{max\_actions}
(20-30), suggesting that comprehensive coverage requires substantial test
generation resources. The Random strategy dominates the best configurations,
with four out of five top positions, achieving coverage rates between 97.1\%
and 100\%.

\begin{table}[h!]
\centering
\caption{Top 5 Configurations by Total Coverage}
\label{tab:best_configurations}
\begin{tabular}{|l|c|c|c|c|c|}
\hline
\textbf{Strategy} & \textbf{Max Tests} & \textbf{Max Actions} & \textbf{State (\%)} & \textbf{Action (\%)} & \textbf{Total (\%)} \\
\hline
Random & 200 & 30 & 100.0 & 100.0 & 100.0 \\
\hline
Random & 200 & 20 & 100.0 & 94.1 & 97.1 \\
\hline
Random & 100 & 30 & 100.0 & 94.1 & 97.1 \\
\hline
Random & 200 & 15 & 100.0 & 94.1 & 97.1 \\
\hline
DepthFirst & 200 & 5 & 87.5 & 88.2 & 87.9 \\
\hline
\end{tabular}
\end{table}

Interestingly, the single Depth-First configuration in the top five uses a
minimal \texttt{max\_actions} value of 5, suggesting that this strategy performs
best with shorter, focused test sequences rather than extended exploration paths.

\subsubsection{Parameter Effect Analysis}
Tables~\ref{tab:max_tests_effect} and \ref{tab:max_actions_effect} quantify how
individual parameters influence coverage outcomes. The \texttt{max\_tests}
analysis reveals dramatic differences in strategy sensitivity. The Random
strategy shows remarkable improvement from 67.6\% to 98.5\% coverage as
\texttt{max\_tests} increases from 10 to 200, demonstrating strong scalability
with additional test generation resources.

\begin{table}[h!]
\centering
\caption{Effect of Max Tests on Total Coverage (max\_actions=20)}
\label{tab:max_tests_effect}
\begin{tabular}{|c|c|c|}
\hline
\textbf{Max Tests} & \textbf{DepthFirst} & \textbf{Random} \\
\hline
10 & 67.6 & 67.6 \\
\hline
25 & 70.6 & 83.8 \\
\hline
50 & 73.5 & 87.5 \\
\hline
100 & 74.3 & 94.9 \\
\hline
200 & 74.3 & 98.5 \\
\hline
\end{tabular}
\end{table}

In contrast, the Depth-First strategy plateaus around 74\% coverage regardless
of the test limit increase, highlighting its inherent limitation in exploring
diverse execution paths. This plateau effect occurs because deterministic
path-following quickly exhausts unique sequences in a finite state machine.

\begin{table}[h!]
\centering
\caption{Effect of Max Actions on Total Coverage (max\_tests=50)}
\label{tab:max_actions_effect}
\begin{tabular}{|c|c|c|}
\hline
\textbf{Max Actions} & \textbf{DepthFirst} & \textbf{Random} \\
\hline
5 & 83.8 & 56.6 \\
\hline
10 & 67.6 & 77.2 \\
\hline
15 & 67.6 & 77.9 \\
\hline
20 & 67.6 & 87.5 \\
\hline
30 & 67.6 & 91.5 \\
\hline
\end{tabular}
\end{table}

The \texttt{max\_actions} analysis reveals counterintuitive behaviour in the
Depth-First strategy, where performance actually decreases from 83.8\% to 67.6\%
as sequence length increases from 5 to 10 actions, then remains stable. This
suggests that longer sequences in deterministic exploration may lead to
repetitive paths that miss alternative state transitions. The Random strategy
shows consistent improvement with longer sequences, reaching 91.5\% coverage
with 30-action sequences, confirming the value of extended exploration for
comprehensive testing.

\subsubsection{Strategic Implications}
The analysis reveals clear strategic implications for test generation parameter
selection. For maximum coverage in exploratory testing scenarios, the Random
strategy with high \texttt{max\_tests} (100-200) and \texttt{max\_actions}
(20-30) values provides optimal results. However, this comes at the cost of
increased execution time and potential test redundancy.

For scenarios prioritising predictability and test suite maintainability, the
Depth-First strategy with moderate parameters offers reliable, consistent
coverage around 75\%. The counterintuitive relationship between \texttt{max\_actions}
and Depth-First performance suggests that shorter, focused test sequences may be
more effective for systematic exploration.

The All-Pairs strategy's limited applicability to rule-constrained systems
represents a significant limitation for complex state machines. Future work
should investigate hybrid approaches that combine the consistency of Depth-First
with the exploration capability of Random strategies, potentially using
adaptive parameter selection based on real-time coverage feedback.
